\section{确定性的分析与估值}
\label{sec:valuation}

真值表生成之后，所有数值计算都由确定性 Python 完成。模型在这一层的作用只有两个：起草估值假设（每一项写出依据并标为“假设”），以及之后在文稿中解读产物。

\subsection{财报分析}
同比与环比变动统一定义为
\begin{equation}
  g = \frac{x_t - x_{t-1}}{|x_{t-1}|}\times 100,
\end{equation}
分母取绝对值，以符合 A 股对基期为负时的惯例。异常信号由固定规则产生，并附严重度与证据定位：非经常性损益占归母净利润比例过高（S1）；归母净利润增速显著高于扣非增速（S2）；净现比偏低（S3）或净利润增长而经营现金流下降（S4）；应收账款、存货增速显著快于收入（S5、S6）；毛利率大幅变动（S7）等。利润变动归因桥把归母净利润的变动拆到毛利、各项费用、投资收益、减值、所得税与少数股东损益上。业绩对比模块给出实际值与一致预期的超预期或不及预期判断、与业绩预告区间中值的比较，以及每股收益质量筛查。

\subsection{估值链：分部预测驱动三表，三表驱动 DCF}
\label{sec:val-chain}

\begin{figure}[t]
  \centering
  \figslot{fig3_valuation}{5.2cm}
  \caption{估值链。经营数据驱动分部预测（销量 × 单价 → 收入，毛利率与费用率），分部预测驱动三表联动模型，DCF 的全部情景向量取自三表模型的同一情景；企业价值经逐项列示的股权价值桥得到每股价值，评级由程序按上行空间给出。虚线表示逐期一致性核对。}
  \label{fig:valuation}
\end{figure}

图 \ref{fig:valuation} 给出估值链。设公司有分部 $s=1,\dots,S$，预测期 $t=1,\dots,N$。

\paragraph{分部预测。}分部收入由量价驱动，
\begin{equation}
  \mathrm{Rev}_{s,t} = Q_{s,t}\cdot P_{s,t},\qquad Q_{s,t}=Q_{s,t-1}(1+g^{Q}_{s,t}),\quad P_{s,t}=P_{s,t-1}(1+g^{P}_{s,t}),
\end{equation}
分部成本由分部毛利率决定，再汇总为公司收入增速与综合毛利率。基期分部收入加其他收入必须等于报表营业收入（评分条目 L1-02）。没有量价数据的行业（例如航运），可以直接给出合并收入增速与毛利率路径，并在分析框架中说明原因。

\paragraph{三表联动。}三表模型按收入、毛利率、费用率、营运资本周转、资本开支与折旧、分红率等驱动预测利润表、资产负债表与现金流量表。经营费用作为轧差项（毛利 $-$ EBIT $-$ 折旧摊销），保证历史期 EBIT 与 DCF 口径一致。模型逐期检查资产负债表平衡、现金流与现金余额衔接以及各项滚动（固定资产、债务、权益）。这里的 EBIT 定义为
\begin{equation}
  \mathrm{EBIT} = \text{利润总额} + \text{财务费用} - \text{投资收益} - \text{公允价值变动收益},
\end{equation}
即剔除投资类收益；被剔除收益对应的投资类资产在股权价值桥中作为非经营性资产加回。

\paragraph{FCFF 折现。}DCF 的收入、EBIT、折旧摊销、资本开支与营运资本\textbf{逐期取自三表模型的同一情景}（评分条目 L0-03 逐期核对其偏差），假设文件中的利润率等字段不再单独生效。
\begin{align}
  \mathrm{FCFF}_t &= \mathrm{EBIT}_t(1-\tau_t) + \mathrm{DA}_t - \mathrm{Capex}_t - \Delta\mathrm{NWC}_t,\\
  \mathrm{WACC} &= w_E K_E + w_D K_D(1-\tau) + w_P K_P,\qquad K_E = r_f + \beta\cdot\mathrm{ERP} + \text{规模、特定与国别溢价},\\
  \mathrm{TV} &= \frac{\mathrm{FCFF}_N(1+g_\infty)}{\mathrm{WACC}-g_\infty}\quad\text{或}\quad \mathrm{TV}=\mathrm{EBITDA}_N\times \text{退出倍数},\\
  \mathrm{EV} &= \sum_{t=1}^{N}\frac{\mathrm{FCFF}_t}{(1+\mathrm{WACC})^{t-0.5}} + \frac{\mathrm{TV}}{(1+\mathrm{WACC})^{N}}.
\end{align}
现金流按年中折现，终值按第 $N$ 年末折现。股权价值桥逐项列示：
\begin{equation}
\begin{split}
  \text{股权价值} = {}& \mathrm{EV} + \text{现金} + \text{非经营性资产} + \text{联营投资}\\
  & - \text{有息负债} - \text{租赁负债} - \text{少数股东权益} - \cdots,
\end{split}
\end{equation}
每股价值 $=$ 股权价值 $/$ 稀释后总股本。非经营性资产须逐项核对年报原文。三个情景（悲观、基准、乐观）各自对应三表模型的一组驱动；敏感性矩阵覆盖 WACC × 永续增长率等组合，程序检查其单调方向（WACC 上升价值下降、永续增长上升价值上升）。

\paragraph{可比公司与交叉验证。}可比公司估值逐家计算企业价值桥与倍数，处理负值与无意义倍数，按四分位距剔除离群值，给出 P25～P75 的估值区间；另输出利润率对比页，用于判断倍数差异是否有基本面支撑。

\paragraph{反向 DCF：股价隐含了什么。}正向 DCF 由假设推出价值；反向 DCF 固定其余假设为基准情景，只动一个变量 $\theta$（预测期收入增速整体平移、EBIT 利润率整体平移、WACC 平移、永续增长率或退出倍数），用网格扫描加二分法求解
\begin{equation}
  V_{\text{每股}}(\theta) = P_{\text{参考}},
\end{equation}
得到参考股价隐含的取值，并与基准假设和五年历史区间对照分档。各变量单独求解而非联合求解，且只作用于 DCF 层、不回推三表模型；这些口径简化写入产物说明。反向 DCF 为研报的“预期差”一节提供了可以复算的锚点。

\paragraph{程序评级。}评级不由模型给出，而由程序按基准每股价值相对参考股价的上行空间 $u = V_{\text{base}}/P_{\text{参考}} - 1$ 决定：
\begin{equation}
  \text{评级} = \begin{cases}
    \text{买入} & u \ge 20\%\\
    \text{增持} & 5\% \le u < 20\%\\
    \text{中性} & -5\% \le u < 5\%\\
    \text{减持} & -20\% \le u < -5\%\\
    \text{卖出} & u < -20\%
  \end{cases}
\end{equation}
研报只引用这一结果。这样评级与估值之间不会出现模型“看好公司所以给买入”的漂移（评分条目 L0-05）。

\subsection{估值模型工作簿}
估值结果同时输出为按投行体例排版的工作簿：输入单元格蓝色、公式黑色、跨表引用绿色；分部预测 → 三表假设（含情景开关）→ 三表全公式联动与勾稽检查 → DCF，历史财务、分季度与经营数据各占一页。关键计算以公式而非数值写入，使用者可以改假设看结果。工作簿审计模块检查跨期公式孤岛（判失败）、公式内写死的数字与易变函数（提示）。

\subsection{与参照引擎的一致性}
三表、情景表、业绩对比、DCF 与可比公司引擎的计算口径参照了一个外部投研插件（细节与许可证说明见仓库 \file{THIRD\_PARTY.md}），但实现全部自研、仓库中不包含其代码。一致性由两类测试保证：\file{tests/test\_golden.py} 要求自研引擎在样例上的产物与改造前原引擎的快照完全相等；在本地存在原插件时，\file{tests/test\_parity\_reference.py} 用样例、边界值与数百组随机输入逐项对照（三表浮点逐位相等，DCF 相对误差小于 $10^{-9}$）。有意修改口径时同步更新快照，并在变更说明中写明原因。
